\documentclass{amia}
\usepackage{amsmath}

\setlength{\bibsep}{0pt}

\begin{document}

\title{KM Curve Analyzer: LLM-Assisted Kaplan-Meier Digitization\\
       and Indirect Treatment Comparison}

\author{Yongxian Zhang}

\institutes{
    Department of Population Health Sciences, Weill Cornell Medicine, New York, NY 10065, US
}

\maketitle

\section*{Abstract}

\textit{Digitizing Kaplan-Meier (KM) survival curves from published figures is
central to evidence synthesis and indirect treatment comparison, yet the process
remains largely manual and time-consuming. We present KM Curve Analyzer, a
web application that integrates large language model (LLM)-based figure
digitization with standard survival statistics. Given a KM figure --- uploaded
directly or retrieved via a PubMed query --- the system extracts
time--survival coordinates, reconstructs individual patient data (IPD) via the
Guyot algorithm, and computes log-rank p-values and hazard ratios with 95\%
confidence intervals. A Bucher indirect comparison module further estimates
the relative effect between two treatments lacking a direct head-to-head
trial. The pipeline supports OpenAI, Anthropic, and Google LLM providers and
handles figures with more than two treatment groups through an interactive
selection interface. Applied to three open-access oncology RCT figures, extracted hazard
ratios agreed with published values with absolute log-scale errors
of 0.00--0.25, with the closest agreement on figures containing
clearly visible at-risk tables.}

\section*{Introduction}

Kaplan-Meier curves are the standard visualization of time-to-event outcomes
in randomized clinical trials (RCTs)\cite{kaplan1958}. Systematic reviewers
and health technology assessors frequently need to recover individual patient
data (IPD) from these figures when original data are unavailable, for example
in network meta-analysis or model-based cost-effectiveness analysis.
Conventional approaches rely on manual digitization tools, which are
time-consuming and operator-dependent\cite{tierney2007}.

Beyond single-trial analysis, indirect treatment comparison (ITC) is
increasingly used to estimate the relative effect of two treatments that have
never been directly compared in a single trial. The Bucher method provides a
frequentist estimator for the indirect hazard ratio when two trials share a
common comparator arm\cite{bucher1997}.

We developed KM Curve Analyzer to automate the full extraction-to-inference
pipeline. A user provides a KM figure (or a PubMed keyword query), and the
system returns log-rank test statistics, a hazard ratio with confidence
interval, and optionally an indirect comparison result --- all without manual
data entry.

\section*{Methods}

\paragraph{LLM-based curve extraction.}
A multimodal prompt is submitted to the selected LLM provider (GPT-4o,
Claude Sonnet, or Gemini Flash). The model returns a structured JSON object
containing, for each group, a list of (time, survival) coordinate pairs,
at-risk table entries, censoring mark positions, and axis metadata including
scale and units. Three providers are supported to allow fallback when
any single provider is unavailable.

\paragraph{IPD reconstruction.}
Individual patient data are reconstructed from the extracted curve
coordinates and at-risk counts using the algorithm of Guyot et
al.\cite{guyot2012}. The algorithm inverts the Kaplan-Meier estimator to
impute event and censoring times consistent with the observed curve shape.
The resulting pseudo-IPD enables standard survival analysis on
digitized data.

\paragraph{Log-rank test and hazard ratio.}
A log-rank test\cite{mantel1966} is performed on the reconstructed IPD.
The hazard ratio and 95\% confidence interval are estimated using a Cox
proportional hazards model\cite{cox1972} (lifelines library); the
Mantel--Haenszel estimator is used as a fallback when Cox model convergence
fails.

\paragraph{Bucher indirect comparison.}
Given two direct comparisons A\,vs\,B and B\,vs\,C sharing a common
comparator B, the indirect log hazard ratio for A\,vs\,C is:
\[
  \ln\widehat{HR}_{AC} \;=\; \ln\widehat{HR}_{AB} \;-\; \ln\widehat{HR}_{BC},
  \qquad
  \mathrm{SE}^2(\ln\widehat{HR}_{AC}) \;=\;
      \mathrm{SE}^2_{AB} + \mathrm{SE}^2_{BC}
\]
A two-sided z-test yields the p-value and 95\% CI\cite{bucher1997}.
Standard Bucher-method assumptions (homogeneity of effect modifiers across
trials) are listed explicitly alongside the result.

\section*{System Design}

The backend is a FastAPI application exposing 13 REST endpoints covering LLM
extraction, IPD reconstruction, log-rank analysis, PubMed search, and PMC
figure retrieval. The frontend is built with Next.js 14 and Tailwind CSS,
providing a browser-based interface with no installation required.

Users may supply figures in two ways: direct image upload, or a PubMed
keyword search that automatically retrieves open-access PMC figures via NCBI
E-utilities and HTML scraping of article pages. For articles behind a
paywall, the interface displays a notification and presents a manual upload
fallback. When a figure contains more than two treatment groups, a two-step
flow extracts all group names first, then presents an interactive selector
before computing statistics, preventing silent truncation errors. API keys
for each LLM provider are stored in browser localStorage.

\section*{Results and Discussion}

\paragraph{Single-figure validation.}
The system was applied to KM figures from three published oncology RCTs
(Table~\ref{tab:validation}). For each figure, the extracted hazard ratio
was compared to the value reported in the source publication. A systematic
direction inversion was observed across all experiments: the Cox model
returns HR(B/A) where group A is the first group listed by the LLM, which
is often the treatment arm. Published papers typically express
HR(treatment/control). The corrected HR (1/raw) therefore corresponds to
the published convention.

\begin{table}[H]
\begin{center}
\begin{tabular}{|p{4.2cm}|l|l|l|l|}
\hline
Trial (endpoint) & Published HR & Extracted HR$^\dagger$ & $|\Delta\ln\text{HR}|$ \\
\hline
CABOSUN\cite{cabosun2017} (PFS) & 0.66 (0.46–0.95) & 0.66 (0.66–0.95) & 0.00 \\
\hline
CheckMate 9ER\cite{checkmate9er2021} (OS) & 0.60 (0.40–0.89) & 0.70 (0.49–1.00) & 0.15 \\
\hline
FRESCO subgroup\cite{fresco2021} (OS) & 0.59 (0.45–0.77) & 0.76 (0.61–0.97) & 0.25 \\
\hline
\end{tabular}
\end{center}
\caption{Single-figure HR validation. $^\dagger$Extracted HR is the
direction-corrected value (1/raw); 95\% CI bounds are also inverted.
$|\Delta\ln\text{HR}|$ measures the absolute log-scale error.}
\label{tab:validation}
\end{table}

The CABOSUN trial yielded exact agreement ($|\Delta\ln\text{HR}|=0.00$),
with patient counts matching the at-risk table (Cabozantinib 79,
Sunitinib 78). The CheckMate 9ER OS figure produced a HR of 0.70 versus
published 0.60 ($|\Delta\ln\text{HR}|=0.15$), with patient counts
approximately matching (extracted 324/328 vs.\ enrolled 323/328).
The FRESCO subgroup figure showed the largest error ($|\Delta\ln\text{HR}|=0.25$),
attributed to a multi-panel layout that caused the LLM to extract
mis-sized at-risk tables.

\paragraph{Multi-panel figure handling.}
A four-panel supplementary figure from the CORRECT trial\cite{correct2015}
(Japanese and non-Japanese subpopulations, OS and PFS) previously caused
group-merging errors, with only 2 groups extracted instead of 8.
After updating the extraction prompt to require panel-prefixed group names,
the system correctly identified all 8 groups (e.g., ``a --- Placebo'',
``b --- Regorafenib''), enabling the user to select any within-panel pair
for analysis.

\paragraph{Indirect comparison.}
As a proof of concept, the Bucher module was applied to two third-line mCRC
trials sharing a placebo/best-supportive-care common comparator: the FRESCO
trial\cite{fresco2021} (Fruquintinib vs.\ Placebo, HR\,1.75, 95\% CI
1.34--2.28) and the CORRECT subpopulation analysis\cite{correct2015}
(Placebo vs.\ Regorafenib, HR\,0.88, 95\% CI 0.35--2.21). The estimated
indirect HR for Fruquintinib vs.\ Regorafenib was 1.53 (95\% CI 0.58--4.00,
$p=0.39$). The wide confidence interval reflects uncertainty propagated from
the imprecise B\,vs\,C input; the workflow itself ran end-to-end from
PubMed search to indirect result without manual data entry.

\paragraph{Limitations.}
Three recurring failure modes were identified: (1) systematic HR direction
inversion due to LLM group-ordering (resolvable by a swap-groups button in
the UI); (2) patient-count over-estimation when at-risk tables are
absent from the figure; and (3) curve shape errors in low-resolution or
multi-panel supplementary figures. Extraction quality is also bounded by
the visual resolution of the published figure and by the exchangeability
assumption required for valid indirect comparison.

\subsection*{Conclusion}

KM Curve Analyzer automates the extraction-to-inference pipeline for
published survival curves. Integration of LLM-based digitization with IPD
reconstruction and indirect comparison methods reduces a multi-hour manual
workflow to minutes. Future work includes systematic validation against a
curated benchmark of digitized KM figures and extension to network
meta-analysis for multi-arm treatment networks.

\makeatletter
\renewcommand{\@biblabel}[1]{\hfill #1.}
\makeatother

\bibliographystyle{vancouver}
\bibliography{report}

\end{document}
